% !TeX root = GeoDynamics.tex
\section{Contact Prediction Processing}

This section documents the contact prediction process. Contact prediction is
the operation that takes one active contact and produces a source-particle
velocity prediction for that contact. The schedule-wise resolution algorithm
then combines these predictions.

Prediction is source-owned. A prediction may read the target particle and the
source-owned contact state, but it returns only a velocity for the source
particle.

\subsection{Particle Contact Prediction Algorithm}

For a source particle \(i\) and target particle \(j\), the particle contact
prediction algorithm is:

\begin{enumerate}
    \item Read the current contact geometry:
    \[
        \mathbf{n}_{ij},\quad A,\quad d .
    \]
    Here \(\mathbf{n}_{ij}\) points from source particle \(i\) to target
    particle \(j\), \(A\) is the current overlap area, and \(d\) is the current
    center distance.

    \item Compute the current relative normal velocity:
    \[
        u_{ij}
        =
        \left(
            \mathbf{v}_j
            -
            \mathbf{v}_i
        \right)
        \cdot
        \mathbf{n}_{ij}.
    \]

    \item Choose the provisional phase:
    \[
        u_{ij} < 0 \Rightarrow \text{compression},
        \qquad
        u_{ij} \ge 0 \Rightarrow \text{rebound}.
    \]

    \item Read the source-owned contact state. If the state contains stored
    zero-velocity geometry, update the phase using the current overlap and
    current relative normal velocity. If there is no zero-velocity geometry and
    the state is not already rebound, the prediction is invalid.

    \item Read the first-contact source and target velocities from the contact
    state:
    \[
        \mathbf{v}_{i}^{0},
        \qquad
        \mathbf{v}_{j}^{0}.
    \]
    If these are missing, the current velocities are used as fallback values.

    \item Construct a two-particle Neo prediction using the current geometry,
    first-contact velocities, masses, radii, and stored zero-velocity geometry.

    \item Compute the Neo report at the current overlap. The report contains
    both compression and rebound velocities:
    \[
        \mathbf{v}_{i,compression}^{\ast},
        \mathbf{v}_{j,compression}^{\ast},
        \mathbf{v}_{i,rebound}^{\ast},
        \mathbf{v}_{j,rebound}^{\ast}.
    \]

    \item Update the diagnostic internal momentum values stored in the contact
    state.

    \item Select the source and target predicted velocities from the report:
    \[
        \text{compression}
        \Rightarrow
        \left(
            \mathbf{v}_{i,compression}^{\ast},
            \mathbf{v}_{j,compression}^{\ast}
        \right),
    \]
    \[
        \text{rebound}
        \Rightarrow
        \left(
            \mathbf{v}_{i,rebound}^{\ast},
            \mathbf{v}_{j,rebound}^{\ast}
        \right).
    \]

    \item If the phase is rebound, apply the outward escape rule to the pair of
    predicted velocities. This enforces a small minimum outward relative normal
    velocity:
    \[
        u_{out} \ge u_{min}.
    \]

    \item If the phase is rebound, apply the internal-contact momentum feedback
    term to the source prediction.

    \item Preserve the source particle's current tangential velocity. Only the
    predicted normal component is transferred into the returned source
    velocity:
    \[
        \mathbf{v}_{i}^{return}
        =
        \mathbf{v}_i
        -
        \left(
            \mathbf{v}_i \cdot \mathbf{n}_{ij}
        \right)
        \mathbf{n}_{ij}
        +
        \left(
            \mathbf{v}_{i}^{\ast} \cdot \mathbf{n}_{ij}
        \right)
        \mathbf{n}_{ij}.
    \]

    \item Return:
    \[
        \mathbf{v}_{i}^{0},
        \qquad
        \mathbf{v}_{i}^{return}.
    \]
\end{enumerate}

\subsection{Wall Contact Prediction Algorithm}

Wall contact prediction uses the same contact prediction machinery with a wall
ghost particle. The wall ghost is fixed and is not integrated.

For source particle \(i\) and wall \(w\), the wall prediction algorithm is:

\begin{enumerate}
    \item Build the wall ghost contact geometry:
    \[
        \mathbf{n}_{w},\quad A,\quad d.
    \]
    The normal \(\mathbf{n}_{w}\) points from the particle toward the wall
    ghost.

    \item Choose the velocity accumulator to process. If no accumulator is
    provided, use the particle's current velocity:
    \[
        \mathbf{v}_{acc} = \mathbf{v}_i .
    \]

    \item Compute the wall normal velocity:
    \[
        u_w
        =
        \mathbf{v}_{acc}
        \cdot
        \mathbf{n}_{w}.
    \]

    \item Choose the provisional wall phase:
    \[
        u_w > 0 \Rightarrow \text{compression},
        \qquad
        u_w \le 0 \Rightarrow \text{rebound}.
    \]

    \item Read the wall contact state and update the phase against the stored
    zero-velocity geometry. The wall update uses:
    \[
        u_{relative} = -2u_w
    \]
    so the wall ghost path follows the same pair-style phase logic.

    \item Call the generic contact prediction operation with the real particle
    as source and the wall ghost as target.

    \item If the phase is rebound and zero-velocity geometry exists, replace
    the generic pair rebound prediction with the wall-specific rebound
    prediction. The wall-specific prediction reflects the source particle
    against the fixed wall:
    \[
        \mathbf{v}_{full}
        =
        \mathbf{v}^{0}
        -
        2
        \left(
            \mathbf{v}^{0}\cdot\mathbf{n}_{w}
        \right)
        \mathbf{n}_{w}.
    \]

    \item Preserve the current tangential component of the accumulator and
    replace only its normal component:
    \[
        \mathbf{v}_{return}
        =
        \mathbf{v}_{acc}
        -
        \left(
            \mathbf{v}_{acc}\cdot\mathbf{n}_{w}
        \right)
        \mathbf{n}_{w}
        +
        \left(
            \mathbf{v}^{\ast}\cdot\mathbf{n}_{w}
        \right)
        \mathbf{n}_{w}.
    \]

    \item Return:
    \[
        \mathbf{v}_{acc},
        \qquad
        \mathbf{v}_{return}.
    \]
\end{enumerate}

\subsection{Phase Update Algorithm}

The prediction process depends on whether the contact is compressing or
rebounding. The phase update algorithm is:

\begin{enumerate}
    \item Classify the current motion mode from relative normal velocity:
    \[
        u < -\epsilon \Rightarrow \text{accumulating},
    \]
    \[
        u > \epsilon \Rightarrow \text{releasing},
    \]
    \[
        |u| \le \epsilon \Rightarrow \text{blocked}.
    \]

    \item Compare the current overlap area \(A\) to the stored zero-velocity
    overlap area \(A_0\). With tolerance \(q\):
    \[
        A \ge A_0(1-q)
        \Rightarrow
        \text{zero reached}.
    \]

    \item If the contact reversed outward before reaching \(A_0\), promote the
    zero-velocity state to the current geometry:
    \[
        A_0 \leftarrow A,
        \qquad
        d_0 \leftarrow d.
    \]

    \item If zero has been reached, set the phase to rebound.

    \item Otherwise, if motion is accumulating, set the phase to compression.

    \item Otherwise, if motion is releasing or blocked, keep or set the phase
    to rebound.
\end{enumerate}

\subsection{Pseudocode}

\begin{verbatim}
particle_contact_prediction(source i, target j):
    contact = particle_contact(i, j)
    if contact does not exist:
        return invalid

    normal, overlap_area, center_distance = contact
    relative_normal_velocity = dot(v[j] - v[i], normal)

    phase =
        compression if relative_normal_velocity < 0
        rebound otherwise

    state = source_owned_contact_state(i, j)
    zero = zero_velocity_geometry(state)

    if zero exists:
        update_phase(state, overlap_area, center_distance,
                     relative_normal_velocity)
        phase = state.phase
    else if state.phase is not rebound:
        return invalid

    prediction =
        generic_contact_prediction(i, j, contact, state, phase)

    return prediction


generic_contact_prediction(source i, target j, contact, state, phase):
    read first-contact source and target velocities
    build temporary two-particle Neo model
    compute current-overlap Neo report
    update diagnostic internal momentum state

    if phase is compression:
        source_velocity = report.source_compression_velocity
        target_velocity = report.target_compression_velocity
    else:
        source_velocity = report.source_rebound_velocity
        target_velocity = report.target_rebound_velocity
        source_velocity, target_velocity =
            apply_outward_escape(source_velocity, target_velocity)
        source_velocity =
            apply_internal_contact_momentum_feedback(source_velocity)

    source_velocity =
        preserve_current_tangent_and_use_predicted_normal(source_velocity)

    return first_contact_source_velocity, source_velocity


wall_contact_prediction(source i, wall w, accumulator velocity):
    contact = wall_ghost_contact(i, w)
    if contact does not exist:
        return invalid

    normal_velocity = dot(accumulator_velocity, wall_normal)
    phase =
        compression if normal_velocity > 0
        rebound otherwise

    update wall phase using relative velocity -2 * normal_velocity

    prediction =
        generic_contact_prediction(i, wall_ghost, contact, state, phase)

    if phase is rebound:
        prediction = wall_specific_rebound_prediction(...)

    prediction =
        preserve_accumulator_tangent_and_use_predicted_normal(prediction)

    return accumulator_velocity, prediction
\end{verbatim}

\subsection{Known Questions}

The prediction process is source-owned, but it still contains several choices
that affect dense-contact behavior:
\begin{itemize}
    \item when zero-velocity geometry should be promoted to the current overlap,
    \item whether outward escape should be applied inside each contact
          prediction or only after contact predictions are combined,
    \item whether internal-contact momentum feedback should be applied before
          or after schedule-wise combination,
    \item whether wall prediction should use pair-style ghost prediction,
          wall-specific prediction, or a single unified expression.
\end{itemize}
